\pdfoutput=1
\documentclass[12pt]{article}

\usepackage{amsmath, amssymb, amsthm}
\usepackage{graphicx}
\usepackage{booktabs}
\usepackage{array}
\usepackage{longtable}
\usepackage{calc}
\usepackage{etoolbox}
\makeatletter
\patchcmd\longtable{\par}{\if@noskipsec\mbox{}\fi\par}{}{}
\makeatother
\IfFileExists{footnotehyper.sty}{\usepackage{footnotehyper}}{\usepackage{footnote}}
\makesavenoteenv{longtable}
\usepackage[colorlinks=true, linkcolor=blue, citecolor=blue, urlcolor=blue]{hyperref}
\usepackage{geometry}
\geometry{margin=1in}

\makeatletter
\newsavebox\pandoc@box
\newcommand*\pandocbounded[1]{%
  \sbox\pandoc@box{#1}%
  \Gscale@div\@tempa{\textheight}{\dimexpr\ht\pandoc@box+\dp\pandoc@box\relax}%
  \Gscale@div\@tempb{\linewidth}{\wd\pandoc@box}%
  \ifdim\@tempb\p@<\@tempa\p@\let\@tempa\@tempb\fi
  \ifdim\@tempa\p@<\p@\scalebox{\@tempa}{\usebox\pandoc@box}%
  \else\usebox{\pandoc@box}%
  \fi
}
\makeatother
\providecommand{\tightlist}{%
  \setlength{\itemsep}{0pt}\setlength{\parskip}{0pt}}

\title{A Minimal Recursive Triadic Framework for Self-Similar Hierarchical Systems}
\author{J. W. Milton\\[4pt]{\small Clarity Coalition}}
\date{\small 3 April 2026 \\ v1.0}

\begin{document}
\maketitle

\noindent\textbf{Keywords:} recursive triadic framework; N-D-C; self-similar hierarchy; irreducibility; minimality; dynamical systems

\begin{abstract}
We argue that the minimum non-trivial structure capable of sustained
recursive self-organization is a three-role partition with a single
emergent state and two functionally distinct auxiliary variables.
Starting from a binary state space \(\{0, 1\}\) and the requirement that
any interaction involve at least one bit of state difference, we show
that the smallest non-degenerate simplex (the triangle) induces exactly
three directed roles: a \emph{negotiation} variable \(N\) that carries
the running emergent state, a \emph{definition/limitation} variable
\(D\) that bounds, and a \emph{contribution/integration} variable \(C\)
that accumulates. The result is a proof that \(m = 2\) is insufficient
and \(m = 3\) is necessary for non-trivial convergence under mild
functional independence conditions (Lemma 4.2). We then demonstrate that
the triadic partition supports a self-similar hierarchical extension:
any triad can serve as a node in a higher-level triad without altering
the three-role grammar, producing unbounded depth from three principles.
This framework supplies a formal substrate for the tholonic ladder
family of recurrences \href{https://tholonia.github.io/the-book/}{Mil24}
and for a class of phase-resolved transparency scoring systems
\href{https://tholonia.github.io/the-book/}{Mil24}, but the present
paper establishes only the structural core: minimality, irreducibility,
and self-similar nesting. No empirical claims are made beyond the
logical and topological argument within the stated axioms.
\end{abstract}

\section{Introduction}\label{introduction}

Self-organizing systems appear across disciplines: physical pattern
formation {[}CH93{]}, biological morphogenesis {[}Tur52{]},
computational models of emergence {[}Lan90, Wol02{]}, and the study of
hierarchical complexity in cognitive and social systems {[}JLM01{]}. A
persistent question across these fields is: what is the \emph{minimum}
set of structural ingredients required for a system to organize itself
recursively, without external tuning at every step?

This paper proposes an answer grounded in the properties of binary state
spaces and simplex topology. The argument proceeds in four steps:

\begin{enumerate}
\def\labelenumi{\arabic{enumi}.}
\tightlist
\item
  \textbf{Binary state space} (\(\S 2\)). The minimal non-trivial state
  space capable of supporting a difference (an interaction) is the
  Boolean set \(\{0, 1\}\). With a single bit, no interaction occurs;
  two bits are necessary to encode a relation {[}Sha48{]}.
\item
  \textbf{Minimum simplex} (\(\S 3\)). The smallest non-degenerate
  simplex is the 2-simplex (triangle), which has three vertices and
  three edges. The directed adjacency structure of a labeled triangle
  yields three oriented axes, each connecting one outer vertex to one
  inner midpoint.
\item
  \textbf{Triadic partition} (\(\S 4\)). A recurrence on \(m\) real
  variables with a distinguished state variable and two functionally
  independent auxiliary variables requires \(m \ge 3\) (Lemma 4.2). The
  three roles correspond to the three axes of the triangular geometry:
  bounding (Definition axis), accumulating (Contribution axis), and
  instantiating (Instantiation axis).
\item
  \textbf{Self-similar nesting} (\(\S 5\)). A triadic node can be
  embedded as a single element in a higher-level triad, reproducing the
  three-role structure at every scale. This yields an infinite hierarchy
  from three principles.
\end{enumerate}

The framework is \emph{formal and structural}. It does not derive
specific dynamical laws or empirical constants. It says: if a system is
to organize itself under these minimal axioms, it must have at least
three distinct functional roles, and those roles compose self-similarly.
This is a constraint on admissible architectures, not a prediction of
particular outputs.

\textbf{Related frameworks.} The tholonic ladder recurrence family
\href{https://tholonia.github.io/the-book/}{Mil24} instantiates this
framework concretely by showing that five classical mathematical
constants (\(\pi/4\), \(\varphi\), \(e\), \(\sqrt{2}\), \(\ln 2\))
emerge as limits of three-variable recurrences whose branches differ
only in initial seeds and traversal rules. The present paper supplies
the foundational argument for \emph{why three variables, and why these
three roles}. The supply-chain transparency scoring method (TVPCI)
\href{https://tholonia.github.io/the-book/}{Mil24} applies the same
triadic partition in a non-dynamical, evidence-scoring context,
demonstrating that the three-role structure generalizes beyond
recurrence. This paper isolates the structural core common to both
applications.

\textbf{Scope.} We establish minimality and self-similarity within the
stated axioms. We do not claim that every three-variable system is
irreducible, that all self-organizing systems reduce to this scaffold,
or that the framework produces novel predictions about specific natural
phenomena. The claim is narrower: given the axioms of binary state space
and simplex topology, \emph{three} is a lower bound, and the same three
roles compose hierarchically.

\begin{center}\rule{0.5\linewidth}{0.5pt}\end{center}

\section{Binary state space}\label{binary-state-space}

A state space is a set \(S\) of values that a variable may take. A
relation on \(S\) is a mapping \(R: S \times S \to S\). For \(R\) to be
non-trivial (i.e., to produce an effect distinguishable from the
identity), there must exist at least one ordered pair
\((a, b) \in S \times S\) with \(a \neq b\) such that
\(R(a, b) \neq a\).

\textbf{Definition 2.1} (Binary state space). The \emph{binary state
space} is \(S_2 = \{0, 1\}\), the smallest set capable of encoding a
difference.

\textbf{Proposition 2.2} (Minimality of binary encoding). \emph{Any
system whose state is captured by a single variable requires
\(|S| \ge 2\) to support a non-trivial interaction between states.
Equivalently, one bit is the minimum informational unit for encoding a
relation {[}Sha48{]}.}

\emph{Proof.} If \(|S| = 1\), then \(S = \{a\}\) and \(R(a, a) = a\) for
any \(R\), so every possible \(R\) is the identity mapping. A
non-trivial \(R\) requires at least two distinct elements. \(\square\)

The binary state space \(\{0, 1\}\) is thus the fundamental ground of
any discrete system capable of change. Every richer state space
(integers, reals, vectors) is a composite built atop this binary
substrate.

\textbf{Corollary 2.3} (Interaction requires at least two bits).
\emph{Two interacting variables, each drawn from \(S_2\), encode
\(2 \times 2 = 4\) possible configurations, the minimum capacity for a
directionally non-trivial relation between distinct entities.}

This corollary sets a hard lower bound on the complexity of any
interaction: at least two binary positions, producing four states, must
be available before a direction can be assigned. A triangle, as we will
see, uses six positions (three outer, three inner), exactly enough to
encode three directed axes.

\begin{figure}
\centering
\pandocbounded{\includegraphics[width=\linewidth,keepaspectratio]{figures/3_binary-to-triad.png}}
\caption{Figure 1. From binary state space to triadic structure.}
\end{figure}

\emph{Figure 1. A three-panel schematic diagram. The left panel shows a
single binary digit \(\{0,1\}\). The center panel shows a 2x2 lattice of
four configurations arising from two interacting bits. The right panel
shows a labeled triangle with three outer vertices and three inner
midpoints, with three directed arrows connecting each outer vertex to a
distinct inner vertex.}

\begin{center}\rule{0.5\linewidth}{0.5pt}\end{center}

\section{The minimum non-degenerate simplex}\label{the-minimum-non-degenerate-simplex}

A simplex is the generalization of a triangle to arbitrary dimension.
The \(k\)-simplex is the convex hull of \(k+1\) affinely independent
points. The 0-simplex is a point, the 1-simplex is a segment, and the
2-simplex is a triangle.

\textbf{Definition 3.1} (Non-degenerate simplex). A \(k\)-simplex is
\emph{non-degenerate} if its \(k+1\) vertices are affinely independent;
equivalently, if it has non-zero \(k\)-dimensional volume.

\textbf{Proposition 3.2} (The triangle is the minimum non-degenerate
simplex). \emph{The 0-simplex and 1-simplex are degenerate in the sense
that they cannot encode a directed relation between more than two
elements. The 2-simplex (triangle) is the smallest simplex with three
vertices, three edges, and an interior, enabling three distinct directed
axes without self-intersection.}

\emph{Proof.} The 0-simplex (point) has no distinct elements to relate.
The 1-simplex (segment) has two endpoints and one edge; it supports one
directed relation (from \(A\) to \(B\), or the reverse) but cannot
support two orthogonal directed relations without self-intersection in
the plane. The 2-simplex has three vertices \(\{V_1, V_2, V_3\}\) and
three edges \(\{E_{12}, E_{23}, E_{31}\}\). By labeling the vertices and
orienting each edge, three distinct directed axes are obtained. No
structure with fewer than three vertices supports three independent
directed relations without embedding them in a higher-dimensional space.
\(\square\)

\textbf{Definition 3.3} (Inner midpoint vertices). For a labeled
equilateral triangle with vertices \(A, B, C\) (outer vertices), define
three \emph{inner midpoint vertices} \(a, b, c\) as the midpoints of
edges \(BC, CA, AB\) respectively. The triangle then has six labeled
points: three outer and three inner.

The six-point configuration (outer \(A, B, C\); inner \(a, b, c\)) is
the smallest labeled structure that supports three directed axes, each
from an outer vertex to a distinct inner vertex, without any axis
sharing both its source and its target with another axis. This is the
geometric model from which the three functional roles emerge.

\begin{figure}
\centering
\pandocbounded{\includegraphics[width=\linewidth,keepaspectratio]{figures/3_minimum-simplex.png}}
\caption{Figure 2. The minimum non-degenerate simplex and its three
directed axes.}
\end{figure}

\emph{Figure 2. A two-panel labeled triangle diagram. The left panel
shows an equilateral triangle with three outer vertices (A, B, C) and
three inner midpoints (a, b, c), all labeled. The right panel shows the
same triangle with three colored directed axes overlaid: Definition axis
(blue, \(A \to a\)), Contribution axis (red, \(B \to b\)), and
Instantiation axis (green, \(C \to c\)). No two axes share the same
source-target pair.}

\textbf{Remark} (Connection to binary encoding). Three directed axes
require six labeled positions (three sources, three targets), which is
the capacity of three binary bits (\(2^3 = 8\) configurations, with 6
used). The 2-simplex with midpoints is thus the geometric minimalization
of a structure supporting three independent directed relations in a
binary substrate.

\begin{center}\rule{0.5\linewidth}{0.5pt}\end{center}

\section{The irreducible three-role partition}\label{the-irreducible-three-role-partition}

The geometric structure of \(\S 3\) induces a partition of
\emph{functional roles} within a recursive system. We formalize this as
a lemma on the minimum number of variables required for non-trivial
convergence.

\subsection{Formal statement}\label{formal-statement}

\textbf{Definition 4.1} (Recurrence with distinguished state). Let
\(\mathcal{R}\) be a recurrence on \(m\) real variables
\(x^{(1)}, \ldots, x^{(m)}\), where \(x^{(1)}\) is the
\emph{distinguished state variable} and the update rule is of the form

\[x^{(1)}_{k+1} = g\!\left(x^{(1)}_k;\, \alpha_k, \beta_k\right)\]

where \(\alpha_k\) and \(\beta_k\) are derived from the remaining
\(m-1\) variables. A \emph{convergent non-trivial} recurrence is one for
which \(\lim_{k \to \infty} x^{(1)}_k = L\) exists and \(L\) is not a
fixed point of a purely autonomous function \(h(x^{(1)}_k)\) with
\(m \le 2\).

\textbf{Lemma 4.2} (Three variables are necessary). \emph{Let
\(\mathcal{R}\) be a convergent non-trivial recurrence as defined above.
Suppose:}

\begin{enumerate}
\def\labelenumi{\arabic{enumi}.}
\tightlist
\item
  \(\alpha_k\) and \(\beta_k\) are \emph{functionally independent} in
  the sense that \(g\) is not expressible as a function of a single
  combined argument \(h(\alpha_k, \beta_k)\) for all admissible
  \((\alpha_k, \beta_k)\).
\item
  The limit \(L\) is not a trivial fixed point attainable with
  \(m \le 2\).
\end{enumerate}

\emph{Then \(m \ge 3\): the recurrence requires at least the state
variable \(x^{(1)}\) and two functionally independent auxiliary
variables.}

\emph{Proof.} Assume \(m = 2\). Then only one auxiliary variable, say
\(y\), is available. The update must be of the form
\(x^{(1)}_{k+1} = g(x^{(1)}_k; y_k)\). But then either (i) \(y_k\) is
constant, in which case \(g\) reduces to a function of \(x^{(1)}_k\)
alone and the limit, if it exists, is a fixed point of that function
(violating the non-triviality condition), or (ii) \(y_k\) evolves
according to its own rule \(y_{k+1} = h(x^{(1)}_k, y_k)\), which pushes
the effective dimensionality to 2 but still links \(\alpha_k\) and
\(\beta_k\) through a single functional chain. Condition (1) (functional
independence) requires that the two auxiliary contributions enter \(g\)
through structurally distinct pathways: one bounding/constraining the
step magnitude, the other contributing/integrating additive or
multiplicative structure. A single auxiliary variable cannot carry both
roles simultaneously with functional independence, because its value
alone determines both the bound and the contribution, collapsing them
into a single degree of freedom. Hence \(m \ge 3\). \(\square\)

\textbf{Remark.} This lemma is deliberately weak in its hypotheses. It
does not claim that \emph{all} three-variable recurrences are
irreducible; it claims that a system with a distinguished state and two
functionally independent auxiliary roles cannot be reduced to fewer than
three variables without losing the structural distinction between
bounding and integrating. The empirical evidence that the five tholonic
ladder branches cannot be expressed as non-trivial 2-variable
recurrences is presented elsewhere
\href{https://tholonia.github.io/the-book/}{Mil24}; Lemma 4.2 provides
the formal lower bound.

\subsection{The three roles}\label{the-three-roles}

The three functional roles correspond to the three axes of the
triangular geometry (\(\S 3\)) and to the three semantic positions in
the tholonic triad \href{https://tholonia.github.io/the-book/}{Mil24}:

\textbf{Definition 4.3} (Triadic role partition). For a three-variable
recurrence \((N, D, C)\):

\begin{itemize}
\tightlist
\item
  \(N\) (\emph{negotiation}): the running state being iteratively
  refined; the emergent quantity.
\item
  \(D\) (\emph{definition/limitation}): the bounding or constraining
  parameter; what limits the state at each step.
\item
  \(C\) (\emph{contribution/integration}): the accumulating or
  synthesizing parameter; what drives growth or combines past states.
\end{itemize}

These roles are assigned by function, not by numerical value. Two
variables may share the same numerical seed (e.g., \(D_0 = C_0 = 2\) for
the \(\sqrt{2}\) branch
\href{https://tholonia.github.io/the-book/}{Mil24}) while remaining
operationally distinct: one bounds, the other integrates.

\textbf{Proposition 4.4} (Role stability under seed degeneracy).
\emph{If \(D_0 = C_0\) numerically but the recurrence rule applies \(D\)
and \(C\) through distinct operations, the roles are not degenerate. The
operational distinction (bounding vs.~integrating) persists even when
the numerical values coincide.}

\emph{Proof.} By construction: \(D\) appears in the recurrence as a
divisor that limits the correction magnitude, while \(C\) appears as a
divisor that averages or synthesizes. The two operations are
functionally independent even when \(D_k = C_k\) for all \(k\), because
removing \(D\) from the recurrence eliminates the bounding term, while
removing \(C\) eliminates the synthesis term. The two deletions produce
different dynamical systems. \(\square\)

This proposition is essential: it prevents the objection that ``\(D\)
and \(C\) are the same number, so the triad is really a two-variable
system.'' The numbers may coincide, but the \emph{roles} do not.

\subsection{Geometric correspondence}\label{geometric-correspondence}

The three roles map onto the three directed axes of the labeled triangle
(\(\S 3\)):

\begin{longtable}[]{@{}
  >{\raggedright\arraybackslash}p{(\linewidth - 6\tabcolsep) * \real{0.1935}}
  >{\raggedright\arraybackslash}p{(\linewidth - 6\tabcolsep) * \real{0.3548}}
  >{\raggedright\arraybackslash}p{(\linewidth - 6\tabcolsep) * \real{0.1935}}
  >{\raggedright\arraybackslash}p{(\linewidth - 6\tabcolsep) * \real{0.2581}}@{}}
\toprule\noalign{}
\begin{minipage}[b]{\linewidth}\raggedright
Axis
\end{minipage} & \begin{minipage}[b]{\linewidth}\raggedright
Direction
\end{minipage} & \begin{minipage}[b]{\linewidth}\raggedright
Role
\end{minipage} & \begin{minipage}[b]{\linewidth}\raggedright
Action
\end{minipage} \\
\midrule\noalign{}
\endhead
\bottomrule\noalign{}
\endlastfoot
Definition & \(A \to a\) & \(D\) & Bounds; constrains; subtracts \\
Contribution & \(B \to b\) & \(C\) & Accumulates; integrates; adds \\
Instantiation & \(C \to c\) & \(N\) (derivative) & Instantiates the
emergent state \\
\end{longtable}

The Instantiation axis supplies the step structure that drives \(N\)
forward; \(N\) itself is the running composite of the interactions
between \(D\) and \(C\). The geometry is not decorative: it supplies the
combinatorial constraints on seed values for the \(\pi/4\) branch
\href{https://tholonia.github.io/the-book/}{Mil24}, where the axis
multipliers \((5, 3, 2)\) uniquely determine the seeds \((1, 3, 5)\) and
the step \(\Delta = 4\).

\begin{figure}
\centering
\pandocbounded{\includegraphics[width=\linewidth,keepaspectratio]{figures/3_three-roles.png}}
\caption{Figure 3. Three functional roles of the triadic partition.}
\end{figure}

\emph{Figure 3. A three-panel functional schematic. Each panel shows the
same triangle geometry (vertices A, B, C; midpoints a, b, c) with one
axis highlighted and the other two faint. Left: Definition axis (blue,
\(A \to a\)) annotated ``D: Bounds; constrains; limits.'' Center:
Contribution axis (red, \(B \to b\)) annotated ``C: Accumulates;
integrates; adds.'' Right: Instantiation axis (green, \(C \to c\))
annotated ``N: Emerges; the negotiated state.'' Roles are functionally
distinct even when \(D = C\) numerically.}

\begin{center}\rule{0.5\linewidth}{0.5pt}\end{center}

\section{Self-similarity of the triadic partition}\label{self-similarity-of-the-triadic-partition}

A structure is self-similar if it reproduces the same pattern at
multiple scales. The triadic partition exhibits self-similarity in a
precise sense: a triad can be embedded as a single element within a
higher-level triad, and the embedding preserves the three-role grammar
at every level.

\subsection{Formal definition}\label{formal-definition}

\textbf{Definition 5.1} (Triadic embedding). Let
\(\mathcal{T}_1 = (N_1, D_1, C_1)\) be a triad. A \emph{triadic
embedding} of \(\mathcal{T}_1\) into a higher-level triad
\(\mathcal{T}_2\) is a mapping
\(\iota: \mathcal{T}_1 \hookrightarrow \mathcal{T}_2\) such that
\(\mathcal{T}_1\) occupies one of the three role positions in
\(\mathcal{T}_2\) while retaining its internal three-role structure.

\textbf{Definition 5.2} (Self-similar triadic hierarchy). A sequence of
triads \(\{\mathcal{T}_\ell\}_{\ell=1}^{L}\) is a \emph{self-similar
triadic hierarchy} if for each \(\ell < L\), there exists a role
\(R_\ell \in \{N, D, C\}\) in \(\mathcal{T}_{\ell+1}\) such that
\(\mathcal{T}_\ell\) is embedded at \(R_\ell\) and the internal
structure of \(\mathcal{T}_\ell\) is a triad \((N, D, C)\) governed by
the same three-role grammar.

\textbf{Proposition 5.3} (Unbounded depth). \emph{Given a triadic root
node with roles \((N, D, C)\), the self-similar triadic hierarchy can be
extended to arbitrary finite depth \(L\) without introducing new
structural primitives.}

\emph{Proof.} By induction on \(\ell\). Base case (\(\ell = 1\)): a
single triad \(\mathcal{T}_1 = (N_1, D_1, C_1)\) satisfies the
definition. Inductive step: assume a hierarchy of depth \(\ell\) exists.
Embed \(\mathcal{T}_\ell\) as, say, the \(N\)-role of a new triad
\(\mathcal{T}_{\ell+1} = (\mathcal{T}_\ell, D_{\ell+1}, C_{\ell+1})\).
The resulting structure is again a triad with three roles. The process
can be repeated for any finite \(\ell\). \(\square\)

\begin{figure}
\centering
\pandocbounded{\includegraphics[width=\linewidth,keepaspectratio]{figures/3_self-similar-nesting.png}}
\caption{Figure 4. Self-similar triadic hierarchy showing three levels
of recursive nesting.}
\end{figure}

\emph{Figure 4. A vertical cascade diagram showing three levels of
self-similar triadic nesting, each level a labeled triangle with the
same three-role grammar. Bottom: a base triad \((N, D, C)\). Middle: the
base triad embedded as the \(N\)-role of a level-2 triad, with new \(D\)
and \(C\) vertices. Top: the level-2 triad embedded as the \(N\)-role of
a level-3 triad. Arrows indicate the embedding direction. The three-role
grammar is preserved at every level without introducing new structural
primitives.}

\subsection{Scale invariance of the three roles}\label{scale-invariance-of-the-three-roles}

An important property of the self-similar hierarchy is \emph{scale
invariance}: the three roles have the same semantics at every level. A
\(D\) at level \(\ell\) is a bounding parameter for the dynamics at that
level, regardless of whether it is internally a triad. A \(C\) at level
\(\ell\) is an accumulating parameter at that level, regardless of its
internal complexity.

This scale invariance is what distinguishes the triadic framework from
generic hierarchical decomposition methods (e.g., hierarchical
clustering, multi-level optimization): the semantics of the three roles
are invariant under embedding. In a generic hierarchy, a node at level
\(\ell\) may have a different functional meaning than a node at level
\(\ell+1\). In the triadic hierarchy, every node at every level is a
\((N, D, C)\) triple, with the same three functional roles.

\textbf{Observation 5.4} (Grammatical closure). \emph{The set of all
finite self-similar triadic hierarchies is closed under composition:
applying the embedding operation to any two hierarchies produces another
hierarchy in the same class.}

This closure property supports the claim that the three-role partition
is not merely sufficient but \emph{complete} for this class: no
additional primitive is needed to extend the framework to arbitrary
depth.

\begin{center}\rule{0.5\linewidth}{0.5pt}\end{center}

\section{Consequences and structural implications}\label{consequences-and-structural-implications}

\subsection{Why not two? Why not four?}\label{why-not-two-why-not-four}

The question of whether \(m = 2\) or \(m = 4\) could serve equally well
has natural answers from the geometry and the irreducibility lemma.

\textbf{Two variables.} Lemma 4.2 establishes that two variables cannot
carry both functionally independent roles (bounding and integrating)
simultaneously. A two-variable system either lacks a bounding term
(leading to divergence) or lacks an integrating term (leading to
stasis). Two variables can capture \emph{binary opposition} (push/pull,
positive/negative) but not the directed emergence of a third state from
the interaction, which requires a place for the emergent quantity to
reside.

\textbf{Four variables.} Adding a fourth variable produces redundancy,
not additional expressive power, at the structural level. A fourth
variable can always be absorbed into one of the three roles without loss
of generality, because the update to the state variable \(N\) depends
only on the bounding contribution (from \(D\), or any composition of
\(D\)-like variables) and the integrating contribution (from \(C\), or
any composition of \(C\)-like variables). If \(m = 4\), the fourth
variable is either (a) a second bounding parameter, which can be
combined with \(D\) into a single effective bound, (b) a second
integrating parameter, which can be combined with \(C\), or (c) a second
state variable, which introduces a new degree of freedom but does not
change the triadic logic at the interface between the state and its
auxiliary parameters. The triadic partition is thus \emph{satisficing}:
it is the smallest structure that meets all requirements, and larger
structures do not add a new functional \emph{kind}.

\textbf{Remark (ternary state spaces).} The argument assumes a binary
state space \(\{0, 1\}\). If the state space were ternary
(\(\{0, 1, 2\}\)), the minimum non-degenerate simplex would remain the
triangle (2-simplex), but the encoding capacity per vertex would
increase. The irreducibility of three roles would persist; only the
combinatorial constraints on seed values would relax. The binary ground
is therefore the \emph{strongest} constraint: if three roles are
necessary in the binary case, they are necessary in any richer state
space as well.

\subsection{Connection to the tholonic ladder}\label{connection-to-the-tholonic-ladder}

The five tholonic ladder branches
\href{https://tholonia.github.io/the-book/}{Mil24} instantiate the
abstract triadic framework with a specific family of recurrences. The
ladder paper establishes:

\begin{itemize}
\tightlist
\item
  A single three-variable recurrence template \((N_k, D_k, C_k)\).
\item
  Five branches distinguished by initial seeds and traversal rules.
\item
  Three traversal classes (Advancing, Self-redefined, Fixed) mapping
  onto different dynamical regimes.
\item
  Seven proved structural properties (diagonal invariance, swap
  symmetry, perfect-square denominators, offset consistency, seed
  partition, convergence to classical limits, and an envelope bound for
  the \(\pi/4\) recursion).
\item
  Convergence to five classical constants (\(\pi/4\), \(\varphi\),
  \(e\), \(\sqrt{2}\), \(\ln 2\)).
\end{itemize}

The present paper supplies the \emph{architectural justification} for
that family: the three-variable template is not an arbitrary choice but
the minimum required by the logic of binary state space, simplex
topology, and functional independence. The ladder paper demonstrates
that the template is \emph{sufficient} to unify five disparate classical
limits; this paper argues that it is \emph{necessary} for any convergent
non-trivial recurrence meeting the stated axioms.

\subsection{Connection to supply-chain transparency scoring}\label{connection-to-supply-chain-transparency-scoring}

The TVPCI framework \href{https://tholonia.github.io/the-book/}{Mil24}
applies the same triadic partition in an evidence-scoring context rather
than a recurrence context. A transparency score for a supply-chain
entity is computed from:

\begin{itemize}
\tightlist
\item
  \(N\): the \emph{negotiated score}, the emergent trustworthiness
  metric.
\item
  \(D\): \emph{definitional evidence}, constraining what counts as
  transparency (documentation, audit trails, certifications).
\item
  \(C\): \emph{contributory evidence}, integrating new information over
  time (incident reports, stakeholder submissions, continuous monitoring
  data).
\end{itemize}

The structural claim is identical: a transparency metric that is both
bounded (by definitional standards) and accumulated (by incoming
evidence) requires at least three variables. Two-variable scoring
systems (e.g., a simple compliant/non-compliant binary, or a single
aggregate score without a bounding reference) cannot simultaneously
track \emph{what counts} and \emph{what has been observed} as
functionally independent quantities.

This demonstrates that the triadic partition is not specific to
recurrence dynamics; it generalizes to any system in which a state is
refined through the interaction of a constraining definition and a
contributing accumulation.

\begin{center}\rule{0.5\linewidth}{0.5pt}\end{center}

\section{Related work}\label{related-work}

\textbf{Minimal models of emergence.} The question ``what is the minimal
system that displays property \(P\)?'' has a long history. Wolfram's
classification of cellular automata {[}Wol02{]} identifies Class 4 rules
(e.g., Rule 110) as the minimal computational universality class within
a 1D binary CA, requiring at least 3-cell neighborhoods. Gardner's Game
of Life {[}Gar70{]} demonstrates emergent complexity from a 2D binary CA
with simple local rules. These examples require many more than three
variables at the micro-level; the triadic framework operates at the
macro-level of \emph{functional description}, where three roles suffice
to describe the dynamics regardless of the micro-implementation.

\textbf{Simplex and category theory.} The 2-simplex as the minimal
non-trivial shape appears in category theory as the basic shape for
composition: a commuting triangle represents the composition axiom
\(f \circ g = h\) {[}ML98{]}. Higher category theory generalizes this to
\(n\)-simplex structures, but the foundational observation remains: the
triangle is the minimum shape encoding directed composition. The triadic
framework repurposes this observation for dynamical systems rather than
morphism composition.

\textbf{Ternary logic and three-valued systems.} Three-valued logics
(true, false, unknown) trace back to \{\L\}ukasiewicz {[}\{\L\}uk20{]}
and have been applied in database theory, circuit design, and
philosophical logic. The triadic framework is not a ternary logic: the
three roles are \emph{functional} (bounding, integrating, emerging), not
\emph{truth-valued}. The connection is structural rather than semantic:
both frameworks recognize that a binary distinction (true/false,
bound/contribute) must be supplemented by a third position for closure.

\textbf{Hierarchical and multi-scale systems.} The literature on
hierarchical systems is vast {[}Sim62, Koe69, Sal12{]}. The distinctive
feature of the triadic hierarchy (\(\S 5\)) is not the existence of
levels but the \emph{grammatical invariance} of the three roles under
embedding. Most hierarchical frameworks allow node semantics to vary
across levels; the triadic framework constrains every node at every
level to be a \((N, D, C)\) triple. This is a stronger claim, and
correspondingly narrower in applicability.

\textbf{Spin networks and twistor theory.} The role-labeled triangle
with inner and outer vertices bears a structural resemblance to spin
networks in loop quantum gravity {[}RS94{]} and to twistor diagrams
{[}Pen67{]}, where combinatorial labeling of simplex edges encodes
quantum states. The resemblance is noted here as a potential direction
for formal connection but is not developed; this paper makes no
quantum-mechanical claims.

\begin{center}\rule{0.5\linewidth}{0.5pt}\end{center}

\section{Discussion}\label{discussion}

\textbf{Generality of the axioms.} The framework rests on three axioms:
(i) binary state space \(\{0, 1\}\) as the minimal substrate, (ii) the
2-simplex as the minimum non-degenerate geometric structure, and (iii)
functional independence of bounding and integrating contributions. These
are not strong axioms. (i) is a consequence of information theory
{[}Sha48{]}; (ii) is a property of Euclidean geometry; (iii) is a
definitional condition on what counts as a \emph{distinct} role. The
result (\(m \ge 3\), self-similar nesting) follows as a structural
deduction rather than an empirical claim.

\textbf{What the framework does not claim.} Several important negative
statements bound the contribution:

\begin{itemize}
\tightlist
\item
  It does not claim that every three-variable system has these three
  roles; only that a system with a distinguished emergent state and two
  functionally independent auxiliary roles corresponds to this
  partition.
\item
  It does not claim that all self-organizing systems reduce to this
  scaffold; only that if a system satisfies the axioms, it admits a
  triadic description.
\item
  It does not derive specific dynamical laws or empirical predictions
  about natural systems; the framework is a constraint on admissible
  architectures, not a generator of observable phenomena.
\item
  It does not assert that the triadic partition is the \emph{only}
  minimal framework; alternative encodings (e.g., quaternary,
  continuous) may exist but would either be reducible to the triadic
  partition or violate the axioms differently.
\end{itemize}

\textbf{Open questions.}

\begin{enumerate}
\def\labelenumi{\arabic{enumi}.}
\tightlist
\item
  \emph{(Classification).} Characterize the set of all three-variable
  recurrences that satisfy the triadic role partition and converge to a
  non-trivial limit. Is this set finite for a bounded seed space?
\item
  \emph{(Embedding completeness).} Is every finite self-similar triadic
  hierarchy representable as a composite of base triads, or do new
  structural types emerge at certain depths?
\item
  \emph{(Continuous limit).} What is the continuous-time analogue of the
  triadic recurrence, and does it correspond to a known class of
  differential equations (e.g., three-variable Lotka-Volterra, Lorenz
  system)?
\item
  \emph{(Observational signatures).} If a physical or biological system
  satisfies the triadic axioms, what measurable signatures would
  distinguish it from a system organized under a different structural
  grammar?
\end{enumerate}

\textbf{Empirical grounding.} The five-constants paper
\href{https://tholonia.github.io/the-book/}{Mil24} provides a concrete
computational instantiation. The TVPCI paper
\href{https://tholonia.github.io/the-book/}{Mil24} provides a
non-dynamical application. Neither claims empirical validation of the
framework as a description of natural phenomena. The next step is to
identify a domain where the three-role partition makes a falsifiable
prediction: either a system that \emph{must} exhibit three functionally
independent roles at a certain level of description, or a system where
the absence of one role predicts a specific type of failure or
instability. Until such a prediction is formulated and tested, the
framework remains a structural taxonomy, albeit one with formal
minimality and self-similarity attached.

\begin{center}\rule{0.5\linewidth}{0.5pt}\end{center}

\section{Conclusion}\label{conclusion}

We have argued that the minimum non-trivial structure capable of
sustained recursive self-organization within a binary state space is a
three-role partition with functionally distinct bounding and integrating
contributions to a distinguished emergent state. The argument derives
from information-theoretic minimality (binary encoding), geometric
minimality (2-simplex), and functional independence (Lemma 4.2),
yielding \(m \ge 3\) as a lower bound. The same three-role grammar
supports unbounded self-similar nesting, producing hierarchical depth
without new primitives.

The framework supplies the architectural foundation for the tholonic
ladder recurrence family
\href{https://tholonia.github.io/the-book/}{Mil24} and for triadic
transparency scoring \href{https://tholonia.github.io/the-book/}{Mil24}.
Its contribution is logical and geometric: minimality, irreducibility,
and hierarchical closure under three axioms. The open work is
classification (which recurrences satisfy the axioms and converge),
embedding completeness (whether new structures emerge at depth), and
empirical connection (can the framework be falsified in a natural
system).

\begin{center}\rule{0.5\linewidth}{0.5pt}\end{center}

\section{References}\label{references}

{[}CH93{]} M. C. Cross and P. C. Hohenberg, ``Pattern formation outside
of equilibrium,'' \emph{Reviews of Modern Physics} 65(3), 1993,
pp.~851--1112.

{[}Gar70{]} M. Gardner, ``The fantastic combinations of John Conway's
new solitaire game `life','' \emph{Scientific American} 223, 1970,
pp.~120--123.

{[}JLM01{]} J. G. Johnson, M. Lebreton, and W. Mansell, ``A
computational model of hierarchical complexity,'' \emph{Behavioral
Development Bulletin} 10(1), 2001, pp.~2--6.

{[}Koe69{]} A. Koestler, \emph{The Ghost in the Machine}, Macmillan,
1969.

{[}Lan90{]} C. G. Langton, ``Computation at the edge of chaos: Phase
transitions and emergent computation,'' \emph{Physica D} 42(1--3), 1990,
pp.~12--37.

{[}\{\L\}uk20{]} J. \{\L\}ukasiewicz, ``O logice trójwartościowej'' (On
three-valued logic), \emph{Ruch Filozoficzny} 5, 1920, pp.~170--171.
English translation in L. Borkowski (ed.), \emph{Selected Works},
North-Holland, 1970.

\href{https://tholonia.github.io/the-book/}{Mil24} J. W. Milton,
\emph{Tholonia: The Existential Mechanics of Awareness}, self-published,
2024. Available at https://tholonia.github.io/the-book/.

{[}ML98{]} S. Mac Lane, \emph{Categories for the Working Mathematician},
2nd ed., Springer, 1998.

{[}Pen67{]} R. Penrose, ``Twistor algebra,'' \emph{Journal of
Mathematical Physics} 8(2), 1967, pp.~345--366.

{[}RS94{]} C. Rovelli and L. Smolin, ``Spin networks and quantum
gravity,'' \emph{Physical Review D} 52(10), 1995, pp.~5743--5759.

{[}Sal12{]} S. Salthe, \emph{Evolving Hierarchical Systems: Their
Structure and Representation}, Columbia University Press, 2012.

{[}Sha48{]} C. E. Shannon, ``A mathematical theory of communication,''
\emph{Bell System Technical Journal} 27(3), 1948, pp.~379--423.

{[}Sim62{]} H. A. Simon, ``The architecture of complexity,''
\emph{Proceedings of the American Philosophical Society} 106(6), 1962,
pp.~467--482.

{[}Tur52{]} A. M. Turing, ``The chemical basis of morphogenesis,''
\emph{Philosophical Transactions of the Royal Society B} 237(641), 1952,
pp.~37--72.

{[}Wol02{]} S. Wolfram, \emph{A New Kind of Science}, Wolfram Media,
2002.

\begin{center}\rule{0.5\linewidth}{0.5pt}\end{center}

\subsection{Appendix A. The 2-simplex as a labeled
configuration}\label{appendix-a.-the-2-simplex-as-a-labeled-configuration}

We provide the explicit vertex and axis assignment for the labeled
triangle used in \(\S\S 3\)--4. The purpose is to make the geometric
origin of the three roles fully transparent, separate from any
particular instantiation.

\textbf{Vertices.} Let the outer vertices be \(A, B, C\) and the inner
midpoints be \(a, b, c\), with: - \(a\) the midpoint of \(BC\) - \(b\)
the midpoint of \(CA\) - \(c\) the midpoint of \(AB\)

\textbf{Axes.} Three directed axes: - \textbf{Definition axis}:
\(A \to a\) (outer to inner, crossing the edge \(BC\)) -
\textbf{Contribution axis}: \(B \to b\) (outer to inner, crossing the
edge \(CA\)) - \textbf{Instantiation axis}: \(C \to c\) (outer to inner,
crossing the edge \(AB\))

\textbf{Binary encoding.} Assigning three bits to encode the six vertex
labels uses 6 of 8 possible configurations:

\begin{longtable}[]{@{}
  >{\raggedright\arraybackslash}p{(\linewidth - 8\tabcolsep) * \real{0.1026}}
  >{\raggedright\arraybackslash}p{(\linewidth - 8\tabcolsep) * \real{0.0897}}
  >{\raggedright\arraybackslash}p{(\linewidth - 8\tabcolsep) * \real{0.2692}}
  >{\raggedright\arraybackslash}p{(\linewidth - 8\tabcolsep) * \real{0.2692}}
  >{\raggedright\arraybackslash}p{(\linewidth - 8\tabcolsep) * \real{0.2692}}@{}}
\toprule\noalign{}
\begin{minipage}[b]{\linewidth}\raggedright
Vertex
\end{minipage} & \begin{minipage}[b]{\linewidth}\raggedright
Label
\end{minipage} & \begin{minipage}[b]{\linewidth}\raggedright
Bit 1 (Outer/Inner)
\end{minipage} & \begin{minipage}[b]{\linewidth}\raggedright
Bit 2 (Axis member)
\end{minipage} & \begin{minipage}[b]{\linewidth}\raggedright
Bit 3 (Role parity)
\end{minipage} \\
\midrule\noalign{}
\endhead
\bottomrule\noalign{}
\endlastfoot
\(A\) & outer & 0 & 1 & 0 (D) \\
\(B\) & outer & 0 & 1 & 1 (C) \\
\(C\) & outer & 0 & 0 & 0 (N) \\
\(a\) & inner & 1 & 1 & 0 (D target) \\
\(b\) & inner & 1 & 1 & 1 (C target) \\
\(c\) & inner & 1 & 0 & 0 (N target) \\
\end{longtable}

The unused two configurations correspond to symmetry-related states not
needed for the three-axis labeling. The three-bit encoding is the
minimal information representation of the geometric structure; each axis
requires one bit for direction (outer \(\to\) inner) and shares bits for
axis identity and role parity.

\begin{center}\rule{0.5\linewidth}{0.5pt}\end{center}

\subsection{\texorpdfstring{Appendix B. Comparison with \(m = 2\) and
\(m = 4\) template
forms}{Appendix B. Comparison with m = 2 and m = 4 template forms}}\label{appendix-b.-comparison-with-m-2-and-m-4-template-forms}

To make the irreducibility argument concrete, we exhibit the forms of
\(m = 2\) and \(m = 4\) recurrences and why they fail or are redundant.

\textbf{Two-variable form.} Let the variables be \((N, \alpha)\). A
recurrence with a bounding or contributing auxiliary role:

\[N_{k+1} = f(N_k; \alpha_k), \qquad \alpha_{k+1} = h(N_k, \alpha_k).\]

If \(f\) is additive, e.g., \(N_{k+1} = N_k \pm 1/\alpha_k\), then
\(\alpha_k\) controls both the sign and the magnitude of the step. The
sign encodes bounding vs.~contributing, but the magnitude is determined
by the same variable. A single \(\alpha\) must serve both as the limiter
and as the integrator, violating functional independence. If \(f\) is
multiplicative, e.g., \(N_{k+1} = N_k \cdot g(\alpha_k)\), the same
issue arises: \(g(\alpha_k)\) determines both the growth factor and the
saturation bound. Two-variable recurrences can converge (e.g., the
logistic map), but they do not separate bounding and contributing into
functionally independent pathways.

\textbf{Four-variable form.} Let the variables be \((N, D_1, D_2, C)\).
The update:

\[N_{k+1} = g(N_k; D_{1,k}, D_{2,k}, C_k).\]

If \(D_{1,k}\) and \(D_{2,k}\) both enter \(g\) as divisors or bounds,
they can be combined into an effective bound
\(D_{\text{eff}} = h(D_1, D_2)\) without changing the functional form of
\(g\). The three-role partition is recovered by identifying
\(D \equiv D_{\text{eff}}\). The fourth variable adds parametric freedom
but not a new functional \emph{kind}: it remains within the
bounding/constraining class. The same argument applies to additional
\(C\)-like variables. Four variables do not break the triadic grammar;
they over-parameterize within it.

\end{document}
